\begin{abstract}
Remote terminal access remains a fundamental interface for operating servers, cloud instances, edge devices, and experimental testbeds. While traditionally driven by human operators, recent automation workflows increasingly exercise remote shells through repeated command execution, output parsing, and interactive terminal updates. These automated high-frequency interactions are highly sensitive to transport-level bottlenecks. SSHv2 remains the de facto protocol for remote shell access, but its TCP transport suffers from head-of-line (HOL) blocking and binds sessions to fixed network addresses, causing severe stalls during automated command execution under network impairments \cite{Scharf2007}. Two modern alternatives target these legacy transport weaknesses: Mosh replaces the TCP byte-stream with a UDP-based terminal state-synchronization model to improve responsiveness under loss \cite{Winstein2012}. While SSH3 redesigns secure remote control over HTTP/3 and QUIC to reduce setup latency and enable stream multiplexing \cite{michel2023}. 



\end{abstract}

\begin{IEEEkeywords}
Remote shell protocols, SSH, Mosh.

\end{IEEEkeywords}